\documentclass[aspectratio=169]{beamer}
\usepackage[english]{babel}
\usepackage[utf8]{inputenc}
\usepackage[T1]{fontenc}
\usepackage{graphicx}
\usepackage{booktabs}
\usepackage{xcolor}

\usetheme{metropolis}
\setbeamercolor{background canvas}{bg=white}
\usecolortheme{default}
\setbeamertemplate{navigation symbols}{}

\setbeamertemplate{footline}
{
  \leavevmode%
  \hbox to \paperwidth{%
    \hfill
    \begin{beamercolorbox}[ht=3ex,dp=2ex,right]{}
       \color{black}
       \insertframenumber{} / \inserttotalframenumber\hspace*{3ex}
    \end{beamercolorbox}%
  }%
  \vskip0pt%
}

\title{Modernizing JADE: Agentic Pipeline Legacy Migration}
\subtitle{Equipping AI Agents with Custom Skills for Codebase Evolution}
\author{Mateusz Jarosz, Sebastian Rydz}
\institute{Warsaw University of Technology}
\date{\today}

\begin{document}

\begin{frame}
  \titlepage
\end{frame}

\begin{frame}{Project Concept and Objective}
\begin{itemize}
  \item \textbf{Problem:} The JADE framework is constrained by outdated Java versions (\texttt{1.5}). Standard LLM chats fail at complex migrations due to context loss and lack of execution tools.
  \item \textbf{Solution:} Build an \textbf{Agentic Pipeline} by empowering an advanced AI Agent (e.g., OpenCode) with a custom suite of skills to collaboratively modernize the codebase.
  \item \textbf{Architecture Focus:} This is a tool-building and methodology task:
  \begin{itemize}
    \item \textbf{The Core:} An existing Chat AI interface acting as the orchestrator.
    \item \textbf{Our Contribution:} Custom CLI tools, testing sandboxes, and structured prompts tailored specifically for JADE and legacy Java.
  \end{itemize}
\end{itemize}
\end{frame}

\begin{frame}{The Agentic Pipeline Architecture}
\begin{itemize}
  \item \textbf{Agentic Paradigm:} Moving beyond simple prompt-response to a continuous, tool-augmented pipeline. 
  \item \textbf{Where is the Multi-Agent?} 
  \begin{itemize}
      \item Currently structured as a \textbf{Router/Orchestrator Agent} using specialized skills.
      \item This lays the architectural groundwork for a true multi-agent system, where skills (Tester, Refactorer, Linter) can easily be decoupled into autonomous sub-agents.
  \end{itemize}
\end{itemize}
\end{frame}

\begin{frame}{The Agentic Migration Methodology}
\begin{itemize}
  \item \textbf{Stepping-Stone Approach:} Execute iterative version jumps (\texttt{1.5 -> 1.6 -> 1.7 -> 1.8 LTS -> newer LTS}).
  \item \textbf{The Agentic Execution Loop:}
  \begin{itemize}
    \item \textbf{Contextualize:} Feed the Agent generated \texttt{CHANGELOG.md} and \texttt{MIGRATION-GUIDE.md}.
    \item \textbf{Instruct:} Prompt the agent to utilize its custom skills on specific JADE modules.
    \item \textbf{Execute \& Iterate:} The agent drafts \texttt{REFACTOR.md}, applies changes, and runs our custom test scripts.
    \item \textbf{Knowledge Capture:} The agent summarizes hurdles into \texttt{LESSONS.md} to improve prompts for the next version jump.
  \end{itemize}
\end{itemize}
\end{frame}

\begin{frame}{Creating Custom Agent Skills}
\begin{itemize}
  \item \textbf{Integration:} Extending the Agent's capabilities via custom bash scripts, CLI wrappers, or MCP integrations.
  \item \textbf{Skill Set Provided to the Agent:}
  \begin{itemize}
    \item \texttt{env-manager}: Spins up isolated Maven/Gradle environments for target Java versions.
    \item \texttt{test-runner}: Triggers JUnit test suites and feeds failure reports directly back into the agent's context.
    \item \texttt{linter-check}: Runs PMD/SonarLint/SpotBugs, allowing the agent to autonomously fix architectural debt.
    \item \texttt{context-fetcher}: Retrieves deprecation lists for specific Java version jumps.
  \end{itemize}
\end{itemize}
\end{frame}

\begin{frame}{Demonstration Scenarios}
\begin{itemize}
  \item \textbf{Scenario 1 -- Skill Verification (\texttt{1.5} to \texttt{1.6}):} The Agent successfully uses \texttt{env-manager} to set up Java 1.6, applies updates, and verifies via \texttt{test-runner}.
  \item \textbf{Scenario 2 -- The \texttt{1.8} LTS Shift:} The developer prompts the modernization of legacy loops. The agent uses linter skills to rewrite inner classes to Lambdas/Streams.
  \item \textbf{Scenario 3 -- Accelerated Legacy:} The \texttt{1.8} to \texttt{1.11} jump requires significantly fewer human interventions because the context from \texttt{LESSONS.md} preempts common errors.
\end{itemize}
\end{frame}

\begin{frame}{Proof of Effectiveness}
\begin{itemize}
  \item \textbf{KPI 1 -- Regression Rate:} Target is \textbf{0} broken core functionalities, verified by the agent's use of JUnit skills.
  \item \textbf{KPI 2 -- Prompt Reduction:} Measuring the decrease in human interventions as the agent learns from previous version jumps.
  \item \textbf{KPI 3 -- Code Quality Score:} Quantifiable decrease in technical debt using integrated static analysis tools.
\end{itemize}
\end{frame}

\begin{frame}[plain]
  \centering
  \Huge \textbf{Thank You!}

  \vspace{1.2em}
  \Large \textbf{Questions?}
\end{frame}

\end{document}